\documentclass{article}
\usepackage[T1]{fontenc}
\usepackage{lmodern}
\usepackage{graphicx}
\usepackage{amsmath,amssymb}
\usepackage[a4paper,margin=2cm]{geometry}
\usepackage[absolute,overlay]{textpos}
\setlength{\TPHorizModule}{1mm}
\setlength{\TPVertModule}{1mm}
\usepackage[indonesian]{babel}
\usepackage{hyperref}
\setlength{\emergencystretch}{2em}
% unit: Halaman 7

\begin{document}

% === Halaman 7 (llm) ===

\begin{itemize}
    \item Nilai entropi berkisar dalam rentang:
    \[ 0 \le H(X) \le \log_2(n) \]
    \item Entropi minimum $= 0$ bit
    \begin{itemize}
        \item $\to$ Tidak ada ketidakpastian
        \item $\to$ Terjadi jika hanya ada satu simbol dengan peluang $1$ (pasti).
    \end{itemize}
    Contoh: $n = 1$, misalkan A dan $p(A) = 1 \to H(X) = -1 \cdot \log_2 1 = 0
    \item Entropi maksimum $= \log_2(n)$ bit
    \begin{itemize}
        \item $\to$ Terjadi jika semua simbol muncul dengan peluang sama, yaitu $1/n$.
        \item $\to$ Ketidakpastian paling tinggi, setiap simbol membawa informasi maksimal.
    \end{itemize}
    Contoh: Empat simbol dengan peluang sama ($0.25$) $\to H(X) = -\sum_{i=1}^{4} 0.25 \log_2 0.25 = 2$ bit
\end{itemize}

\end{document}
